\section{Introduction}\label{sec:intro}

Scheduled macroeconomic releases are among the cleanest information events in financial
markets. Their timing is announced months in advance, a consensus forecast is published before
each release, and the news itself---the difference between the actual figure and the
consensus---arrives at a known second. Prices of index and Treasury futures react within the first
minute \citep{Fleming1999, Balduzzi2001, Andersen2003, Andersen2007}. The speed of that first
reaction is well documented. Much less is known about what happens next: how much of the news is
still unpriced once the first minute has passed, whether the remaining move is predictable from
information that is public at that time, and whether it is large enough to be exploited once
the cost of trading is paid.

These questions matter for two reasons. First, they measure the speed of price discovery in some
of the most liquid markets in the world. If prices fully incorporated the release within a
minute, any later move would be unrelated to the surprise. If they do not, the size and duration
of the remaining drift tell us how long it takes markets to agree on what a number means. Second,
the answer separates two kinds of apparent profitability. A strategy that trades on the surprise
at the release price captures the first-minute jump, which no investor can trade; a strategy that
enters one minute later captures only the drift. Backtests that do not distinguish the two can
overstate what is achievable by a large factor.

We study these questions in four futures markets with different sensitivities to macro
news: the E-mini S\&P 500 (\ES), the E-mini Nasdaq-100 (\NQ), the E-mini Dow (\YM) and the 10-year
Treasury note (\ZN).
We use one-minute prices from July 2010 to June 2026 and the Bloomberg release calendar with
consensus forecasts for 179 U.S.\ release lines. Our design has three features. First, every
trade enters at the close of the first one-minute bar after the release, so the untradeable jump
is excluded by construction. Second, every model parameter---which releases to trade, the sign of
each release's effect, regression slopes and pattern-matching thresholds---is re-estimated each
year using only earlier data, and the last years of the sample are held out. Third, we compare
four forecasting approaches that were developed independently by the authors on the same data,
entry rule and cost grid: (i) the sign of the standardized consensus surprise; (ii) dynamic time
warping (DTW) analogues of the price path just before entry, which ask whether similar paths in
the past were followed by a rise or a fall; (iii) DTW ``event games'' built on longer paths around
the release; and (iv) a machine-learning classifier. The comparison is informative because the
approaches use very different information: the first uses only the published numbers, the second
and third use only prices, and the fourth combines both.

Five findings emerge. First, the news is not fully priced after the first minute in equity-index
futures. After a hot CPI print (a headline surprise above half a standard deviation), \ES\ falls
about 37 basis points (bps) in the first minute and a further 15--18 bps over the next 30 minutes;
\NQ\ falls about 55 bps and then a further 20 bps, and \YM\ about 31 bps and then a further 17
bps. Roughly a quarter of the CPI response arrives
after the first minute. In \ZN\ the post-minute continuation is only 4--5 bps.

Second, the consensus surprise predicts the direction of this drift out of sample. A strategy
that trades the sign of the surprise of CPI alone, entering one minute after the release and
holding for 30 minutes, earns a net Sharpe ratio of about 1.0 in \ES\ and \NQ\ over 2019--2026
after a cost of two ticks plus \$4.50 per round trip. Extending it to about ten release families
selected each year from history raises the trade count tenfold, to about 100 trades a year per
market, with Sharpe ratios of 0.65 (\ES) and 0.71 (\NQ). On the held-out years 2021--2026 at the
lower cost assumption common among practitioners (a quarter tick per side) the Sharpe ratios are
1.09 and 0.75. The same code, applied without change to the E-mini Dow after these results were
known, gives 1.06 on its broad book, and the three equity indices together give 1.07. For the CPI strategy, the direction of the surprise is worth about 13 percentage
points of cumulative return relative to being long in exactly the same 30-minute windows.

Third, price-path analogues do not reliably forecast the direction of the drift. Their rank correlation
with the next 30-minute return is between 0.01 and 0.04 in every market; the sign of their
forecast was positive in 2018--2020 and negative in 2021--2026 in \ES\ and \NQ; and requiring the
analogues to agree with the surprise lowers performance in the held-out years in \ES\ and \NQ\ and
leaves it unchanged in \YM. Run on \YM\ with the same code, the
analogue book loses even before costs. The dispersion of the analogues,
by contrast, is a stable forecast of the size of the move.

Fourth, the drift in \ZN\ is smaller than the cost of trading it. The surprise predicts \ZN's
first-minute jump as well as it predicts \ES's ($R^2$ of about 0.30 at five minutes), and a gross
edge exists after the first minute, but at about 2.7 bps per trade it is below the 3 bps round-trip
cost of one tick. Tick size relative to the size of post-announcement moves acts as a limit to
arbitrage.

Fifth, the entry assumption dominates reported performance. The same classifier positions earn a
Sharpe ratio of 2.1 when entered at the pre-release price and 0.9 when entered one minute later;
about 80\% of the reported profit is the first-minute jump. Across approaches, daily returns are
nearly uncorrelated, and the surprise and DTW books lose in different years, so an equal-risk
combination beats the best single approach: 1.39 on \ES\ and 1.35 on the three equity indices at a
quarter tick, against 1.09, with half the volatility. At realistic costs the high-turnover
pattern-based approaches reduce the combined performance, and the equity portfolio keeps 0.85.

\paragraph{Related literature.}
Our paper contributes to three strands of research. The first studies the high-frequency response
of asset prices to macroeconomic news. \citet{Fleming1999} and \citet{Balduzzi2001} show that
Treasury prices adjust to announcements within minutes, with a burst of volume and volatility;
\citet{Andersen2003, Andersen2007} document that surprises move foreign exchange, bond and stock
prices almost instantly and that the response depends on the state of the economy;
\citet{Gurkaynak2005} show that monetary-policy statements move asset prices beyond the policy
rate itself; and \citet{Savor2013} and \citet{Lucca2015} document return premia on announcement
days and before FOMC meetings. We focus on the part of the response that remains after the first
minute, compare it across equity and Treasury futures, and ask whether it can be exploited net of
costs. The second strand studies how announcement effects depend on market conditions
\citep{Andersen2007}. We find strong in-sample state dependence of the CPI response in equity
futures, concentrated in the untradeable first minute and detectable in real time only from 2022,
and no out-of-sample gain from conditioning on it. The third strand concerns pattern-based
forecasting and the evaluation of trading rules. Dynamic time warping \citep{Sakoe1978,
Berndt1994} is widely used to match time series of different speeds; we provide a clean
out-of-sample test of its directional content and separate it from its information about
volatility. Following \citet{Harvey2016} and \citet{Bailey2014}, we treat the multiplicity of
tested specifications explicitly, and we report random-sign and block-bootstrap inference
\citep{Kunsch1989} alongside heteroskedasticity-robust regressions \citep{MacKinnon1985} and
false-discovery-rate control \citep{Benjamini1995}.

\paragraph{Roadmap.}
Section~\ref{sec:data} describes the institutional setting and data.
Section~\ref{sec:discovery} documents price discovery after releases.
Section~\ref{sec:exploit} sets out the evaluation protocol and the four forecasting approaches.
Section~\ref{sec:results} reports the results, Section~\ref{sec:mechanisms} interprets them,
Section~\ref{sec:robust} presents robustness checks, Section~\ref{sec:future} sets out directions for future research, and Section~\ref{sec:conclusion} concludes.
Appendix~\ref{app:additional} collects additional figures and tables for each stage of the analysis.
